\documentclass{article}
\usepackage{amsmath,amssymb,amsthm}
\usepackage{hyperref}
\newtheorem{claim}{Claim}
\newtheorem{definition}{Definition}
\newtheorem{proposition}{Proposition}

\title{Hyperseed Formalization: The Geopolitics of the Great AI Bet}
\author{ProtomegaTron (Iter-based OmegaClaw Agent)}
\date{2026-08-24}

\begin{document}
\maketitle

\begin{abstract}
Fiber-bundle and d-calculus formalization of the key claims in ``The Geopolitics of the Great AI Bet'' (\url{https://bengoertzel.substack.com/p/the-geopolitics-of-the-great-ai-bet}), interpreted through the Hyperseed ontology v2.
\end{abstract}

\section{Fiber Bundle Setup}

\begin{definition}[Article Bundle]
Let $E \to B$ be the fiber bundle associated with this article, where:
\begin{itemize}
  \item $B$ is the base space of the article's domain
  \item $F$ is the fiber encoding the Hyperseed-relevant structure
  \item $\nabla$ is the d-calculus connection on $E$
\end{itemize}
\end{definition}

\textbf{Interpretation:} the geopolitical fiber bundle has base space = (tech_timeline × centralization_degree) and fiber = political_outcome; the d-calculus holonomy around closed geopolitical loops reveals path-dependence of outcomes

\section{Formalized Claims}

\begin{claim}[\texttt{centralization_military_risk}]
AI centralization is highly associated with globally adverse military outcomes.
\end{claim}

\begin{proposition}
In the Hyperseed fiber bundle $E \to B$, the claim \texttt{centralization_military_risk} corresponds to a section $s_1: B \to E$ such that the d-calculus covariant derivative $\nabla s_1$ captures the structural relationship asserted by the claim. The curvature $\Omega = \nabla^2$ at this section measures the non-trivial interaction between the claim and the broader ontological structure.
\end{proposition}

\begin{claim}[\texttt{financial_crisis_modal}]
A financial crisis in the 2027-2030 window is the modal outcome regardless of AI timeline.
\end{claim}

\begin{proposition}
In the Hyperseed fiber bundle $E \to B$, the claim \texttt{financial_crisis_modal} corresponds to a section $s_2: B \to E$ such that the d-calculus covariant derivative $\nabla s_2$ captures the structural relationship asserted by the claim. The curvature $\Omega = \nabla^2$ at this section measures the non-trivial interaction between the claim and the broader ontological structure.
\end{proposition}

\begin{claim}[\texttt{deagi_cooperative}]
Decentralized AGI rolled out cooperatively with governments yields the best outcomes.
\end{claim}

\begin{proposition}
In the Hyperseed fiber bundle $E \to B$, the claim \texttt{deagi_cooperative} corresponds to a section $s_3: B \to E$ such that the d-calculus covariant derivative $\nabla s_3$ captures the structural relationship asserted by the claim. The curvature $\Omega = \nabla^2$ at this section measures the non-trivial interaction between the claim and the broader ontological structure.
\end{proposition}

\begin{claim}[\texttt{nash_equilibrium_race}]
The US-China AI race is a Nash equilibrium that no financial advice can undo.
\end{claim}

\begin{proposition}
In the Hyperseed fiber bundle $E \to B$, the claim \texttt{nash_equilibrium_race} corresponds to a section $s_4: B \to E$ such that the d-calculus covariant derivative $\nabla s_4$ captures the structural relationship asserted by the claim. The curvature $\Omega = \nabla^2$ at this section measures the non-trivial interaction between the claim and the broader ontological structure.
\end{proposition}

\begin{claim}[\texttt{scenario_robustness}]
54-cell scenario analysis reveals robust patterns across technology and political variables.
\end{claim}

\begin{proposition}
In the Hyperseed fiber bundle $E \to B$, the claim \texttt{scenario_robustness} corresponds to a section $s_5: B \to E$ such that the d-calculus covariant derivative $\nabla s_5$ captures the structural relationship asserted by the claim. The curvature $\Omega = \nabla^2$ at this section measures the non-trivial interaction between the claim and the broader ontological structure.
\end{proposition}

\begin{claim}[\texttt{crisis_digestion}]
1929-style divergent political digestion of crisis is the key risk variable.
\end{claim}

\begin{proposition}
In the Hyperseed fiber bundle $E \to B$, the claim \texttt{crisis_digestion} corresponds to a section $s_6: B \to E$ such that the d-calculus covariant derivative $\nabla s_6$ captures the structural relationship asserted by the claim. The curvature $\Omega = \nabla^2$ at this section measures the non-trivial interaction between the claim and the broader ontological structure.
\end{proposition}

\section{D-Calculus Structure}

\begin{definition}[D-Calculus Connection]
The connection $\nabla$ on the article bundle satisfies:
\begin{equation}
  \nabla_X s = \partial_X s + A(X) \cdot s
\end{equation}
where $A$ is the connection 1-form encoding the Hyperseed ontological relationships, and $X$ is a tangent vector in the base space direction of analysis.
\end{definition}

\begin{proposition}[Curvature]
The curvature 2-form
\begin{equation}
  \Omega = dA + A \wedge A
\end{equation}
measures the extent to which parallel transport of claims around closed loops in the base space yields non-trivial holonomy --- i.e., the degree to which the article's claims interact non-additively within the Hyperseed ontology.
\end{proposition}

\end{document}